\chapter{فك الاقتران والتوافق الفعّال: مبرهنة القيمة المتوسطة لفينوغرادوف}
\personindex{فينوغرادوف، إيفان ماتفييفيتش}
\symbolindex{فك الاقتران والتوافق الفعّال: مبرهنة القيمة المتوسطة لفينوغرادوف}
\label{ch:decoupling-efficient-congruencing-vmvt}

\section*{نطاق الفصل وحالته}

يدرس هذا الفصل متوسطات مجاميع فينوغرادوف، ويشرح كيف قاد مساران مختلفان --- التوافق الفعّال وفك الاقتران \(\ell^2\) للمنحنى اللحظي --- إلى الحد الأمثل في مبرهنة القيمة المتوسطة لفينوغرادوف. لا يعيد الفصل بناء البرهان الكامل لأي من المسارين؛ بل يثبت الهوية العددية والحدود الدنيا البنيوية داخل الموسوعة، ثم ينقل المبرهنة العميقة من مصادرها الأصلية مع حراس صريحة على نطاقها.

\noindent\textbf{حالة الفصل:}
\textenglish{AUTHORED-DRAFT / NON-CITABLE}. فُتحت بوابة التأليف بعد مراجعة مستقلة ناجحة لحزمة ما قبل التأليف، لكن النتائج العشرة تبقى محجوزة وغير قابلة للاستشهاد إلى أن ينجح البناء والمراجعة المستقلة بعد التأليف ويعتمد المالك الفصل صراحة.

\subsection*{حراس النطاق}
\begin{itemize}
  \item نستعمل دائمًا \(e(t)=\exp(2\pi i t)\) والتكامل على \([0,1)^k\).
  \item هوية التعامد والحدان السفليان يثبتان هنا؛ الحد العلوي الأمثل منقول ولا يثبت هنا.
  \item لا يُعرض التوافق الفعّال وفك الاقتران على أنهما برهان واحد أو متكافئان خطوة بخطوة.
\personindex{وولي، تريفور}
  \item نتيجة Wooley لعام 2012 تقدم تأسيسي، والحالة التكعيبية الكاملة تعود إلى عمله اللاحق المنشور سنة 2016.
\personindex{غوث، لاري}
\personindex{بورغين، جان}
\personindex{ديمتر، تشيبرين}
  \item نتيجة Bourgain--Demeter--Guth تخص الدرجات \(k\ge4\)، بينما يغطي التوافق الفعّال المتداخل جميع الدرجات.
\personindex{وارينغ، إدوارد}
  \item تطبيقات وارينغ والمجاميع الأسية جسور تفسيرية إلى الفصلين \ref{ch:circle-method-goldbach-waring} و\ref{ch:exponential-sums-van-der-corput}، وليست إعادة برهنة لنواتهما.
\end{itemize}

\section{مجموع فينوغرادوف والقيمة المتوسطة}
\symbolindex{مجموع فينوغرادوف والقيمة المتوسطة}

\begin{definition}[مجموع فينوغرادوف والقيمة المتوسطة]
\label{def:ch25-vinogradov-mean-value}
\resultid{ANT-DEF-25-01}
\provedhere
لتكن \(k,s\in\mathbb N\) و\(X\ge1\). نضع
\[
e(t)=\exp(2\pi i t),
\]
ونعرف
\[
f_k(\boldsymbol\alpha;X)
=
\sum_{1\le n\le X}
e\!\left(\alpha_1n+\alpha_2n^2+\cdots+\alpha_kn^k\right),
\]
حيث \(\boldsymbol\alpha=(\alpha_1,\ldots,\alpha_k)\in[0,1)^k\). ثم نضع
\[
J_{s,k}(X)
=
\int_{[0,1)^k}
\left|f_k(\boldsymbol\alpha;X)\right|^{2s}
\,d\boldsymbol\alpha.
\]
\end{definition}

هذا المتوسط ليس مجرد معيار تحليلي لحجم مجموع أسي؛ فهو يخزن بالضبط عدد حلول نظام جبري من معادلات القوى المتتالية.

\theoremindex{هوية التعامد وعد حلول نظام فينوغرادوف}
\begin{proposition}[هوية التعامد وعد حلول نظام فينوغرادوف]
\label{prop:ch25-orthogonality-count}
\resultid{ANT-PROP-25-01}
\provedhere
يساوي \(J_{s,k}(X)\) عدد الحلول الصحيحة
\[
1\le x_i,y_i\le X\qquad(1\le i\le s)
\]
للنظام
\[
\sum_{i=1}^{s}x_i^j
=
\sum_{i=1}^{s}y_i^j
\qquad(1\le j\le k).
\]
\end{proposition}

\begin{proof}
نوسع القوة الزوجية على الصورة
\[
|f_k(\boldsymbol\alpha;X)|^{2s}
=
\prod_{i=1}^{s}f_k(\boldsymbol\alpha;X)
\prod_{i=1}^{s}\overline{f_k(\boldsymbol\alpha;X)}.
\]
بعد فتح المجاميع نحصل على مجموع على جميع \(x_i,y_i\in[1,X]\cap\mathbb Z\) للحد
\[
e\!\left(
\sum_{j=1}^{k}\alpha_j
\left(
\sum_{i=1}^{s}x_i^j-
\sum_{i=1}^{s}y_i^j
\right)
\right).
\]
وبما أن
\[
\int_0^1 e(\alpha m)\,d\alpha
=
\begin{cases}
1,&m=0,\\
0,&m\in\mathbb Z\setminus\{0\},
\end{cases}
\]
فإن التكامل في كل متغير \(\alpha_j\) يفرض المعادلة ذات الدرجة \(j\). لذلك لا يبقى إلا الحلول التي تحقق النظام كله، وكل حل يسهم بواحد.
\end{proof}

\section{الوزن الحرج والحدود الدنيا البنيوية}
\symbolindex{الوزن الحرج والحدود الدنيا البنيوية}

\begin{definition}[الوزن الحرج]
\label{def:ch25-critical-weight}
\resultid{ANT-DEF-25-02}
\provedhere
نكتب
\[
K=1+2+\cdots+k=\frac{k(k+1)}2.
\]
هذا هو مجموع درجات معادلات نظام فينوغرادوف، وهو الفقد الطبيعي في أس عدد درجات الحرية.
\end{definition}

\theoremindex{المصدران البنيويان للحد الرئيسي}
\begin{proposition}[المصدران البنيويان للحد الرئيسي]
\label{prop:ch25-structural-lower-bounds}
\resultid{ANT-PROP-25-02}
\provedhere
لكل \(s,k\ge1\) يوجد ثابت موجب يعتمد على \(s,k\) فقط بحيث
\[
J_{s,k}(X)
\gg_{s,k}
X^s+X^{2s-K}.
\]
\end{proposition}

\begin{proof}
أولًا، لكل اختيار مرتب لـ\((x_1,\ldots,x_s)\) نحصل على حل قطري بوضع \(y_i=x_i\). ومن ثم
\[
J_{s,k}(X)\ge \lfloor X\rfloor^s\gg_s X^s.
\]

ثانيًا، نختار ثابتًا صغيرًا \(c_k>0\)، ونقيد التكامل إلى الصندوق
\[
\mathcal B_X=
\left\{
\boldsymbol\alpha:
|\alpha_j|_{\mathbb R/\mathbb Z}
\le c_kX^{-j}
\quad(1\le j\le k)
\right\}.
\]
لكل \(1\le n\le X\) و\(\boldsymbol\alpha\in\mathcal B_X\) لدينا
\[
\left|
\sum_{j=1}^{k}\alpha_jn^j
\right|_{\mathbb R/\mathbb Z}
\le
c_k\sum_{j=1}^{k}\left(\frac nX\right)^j
\le kc_k.
\]
وباختيار \(c_k\) بحيث تقع الأطوار كلها في قوس أقصر من نصف دائرة، نحصل على
\[
|f_k(\boldsymbol\alpha;X)|\gg_k X.
\]
أما حجم الصندوق فهو
\[
\operatorname{meas}(\mathcal B_X)
\asymp_k
\prod_{j=1}^{k}X^{-j}
=X^{-K}.
\]
إذن
\[
J_{s,k}(X)
\ge
\int_{\mathcal B_X}|f_k(\boldsymbol\alpha;X)|^{2s}
\,d\boldsymbol\alpha
\gg_{s,k}X^{2s-K}.
\]
وبجمع الحدين، أو باستخدام أن مجموعهما لا يتجاوز ثابتًا مضروبًا في أكبرهما، نحصل على النتيجة.
\end{proof}

الحدان السابقان يفسران لماذا لا يمكن أن يكون الحد العلوي الصحيح أصغر من
\(X^s+X^{2s-K}\). لكن إثبات حد علوي مطابق، حتى مع خسارة \(X^\varepsilon\)، هو قلب المبرهنة العميقة.

\section{مبرهنة القيمة المتوسطة الرئيسية}
\symbolindex{مبرهنة القيمة المتوسطة الرئيسية}

\theoremindex{المبرهنة الرئيسية للقيمة المتوسطة لفينوغرادوف}
\begin{theorem}[المبرهنة الرئيسية للقيمة المتوسطة لفينوغرادوف]
\label{thm:ch25-main-vmvt}
\resultid{ANT-THM-25-01}
\citedresult[\cite{wooley2016cubic,bourgainDemeterGuth2016vmvt,wooley2019nested}]
لكل \(s,k\in\mathbb N\)، ولكل \(\varepsilon>0\)، لدينا
\[
J_{s,k}(X)
\ll_{s,k,\varepsilon}
X^\varepsilon
\left(
X^s+X^{2s-k(k+1)/2}
\right).
\]
\end{theorem}

مع القضية \ref{prop:ch25-structural-lower-bounds} تصبح هذه الصيغة مثلى حتى عامل \(X^\varepsilon\). عند \(s\le K\) يغلب الحد القطري، وعند \(s\ge K\) يظهر الحد الثاني بوصفه المقياس المستمر أو فوق الحرج. لا يعني هذا أن كل الحلول فوق الحرجة غير قطرية، بل يعني أن الحجم الكلي الأمثل يتغير عند الوزن الحرج.

\subsection*{التسلسل التاريخي المنضبط}

بدأ Wooley التوافق الفعّال في عمله المنشور سنة 2012، وحقق به تقدمًا جذريًا في VMVT وتطبيقاتها \cite{wooley2012efficient}. ثم أثبت الحالة التكعيبية الكاملة \(k=3\) في عمل لاحق نشر سنة 2016 \cite{wooley2016cubic}. وأثبت Bourgain وDemeter وGuth الحدسية للدرجات \(k\ge4\) بواسطة فك الاقتران الحاد للمنحنى اللحظي \cite{bourgainDemeterGuth2016vmvt}. وبعد ذلك قدم Wooley التوافق الفعّال المتداخل، مثبتًا صيغة عامة تشمل جميع الدرجات \cite{wooley2019nested}. هذا التسلسل يمنع نسبة النتيجة العامة كلها إلى ورقة واحدة أو إلى مسار واحد.

\section{مسار التوافق الفعّال}
\symbolindex{مسار التوافق الفعّال}

\begin{remark}[مبدأ التوافق الفعّال]
\label{rem:ch25-efficient-congruencing-principle}
\resultid{ANT-PRIN-25-01}
\citedresult[\cite{wooley2012efficient,wooley2019nested}]
يقسم التوافق الفعّال المتغيرات إلى فئات توافقية modulo قوى أولية، ثم يستعمل نظام معادلات القوى لرفع التطابقات إلى مقاييس أدق. تنتج عملية تكرارية تضغط كتلة المتوسط في عائلات أكثر انتظامًا، مع موازنة دقيقة بين كلفة التقييد والمكسب الناتج من عدم الانحلال. هذا وصف لبنية الطريقة، لا إعادة بناء لسلسلة المتراجحات التقنية الكاملة.
\end{remark}

جوهر الفكرة أن المعادلات
\[
\sum x_i^j\equiv\sum y_i^j\pmod{p^a}
\]
ليست قيودًا مستقلة اعتباطية؛ فعند اختيار متغيرات غير متطابقة modulo \(p\)، تسمح بنية مصفوفة فانديرموند برفع المعلومات التوافقية. في النسخة المتداخلة تصبح المقاييس متعددة الدرجات، وهو ما يمنح الطريقة المرونة اللازمة لإغلاق جميع الدرجات.

\section{مسار فك الاقتران للمنحنى اللحظي}
\symbolindex{مسار فك الاقتران للمنحنى اللحظي}

\theoremindex{فك الاقتران الحاد والجسر إلى VMVT}
\begin{theorem}[فك الاقتران الحاد والجسر إلى VMVT]
\label{thm:ch25-decoupling-bridge}
\resultid{ANT-THM-25-02}
\citedresult[\cite{bourgainDemeter2015decoupling,bourgainDemeterGuth2016vmvt}]
ترتبط مبرهنة VMVT بتقديرات فك اقتران \(\ell^2\) حادة للمنحنى اللحظي
\[
\Gamma(t)=(t,t^2,\ldots,t^k)\subset\mathbb R^k.
\]
تجزئ فترة المعلمة إلى فترات قصيرة، وتفصل مساهمات القطع في معيار \(L^p\) موزون على كرات بالمقياس المناسب. وعند الأس الحرج والتطبيع الصحيحين تقود صيغة فك الاقتران، بعد انتقال متقطع، إلى الحد المطلوب لـ\(J_{s,k}(X)\).
\end{theorem}

لا تكفي مبرهنة فك الاقتران العامة للأسطح ذات الانحناء من دون ضبط هندسة المنحنى اللحظي ومقياس الكرات وطول الفترات. لذلك يجب عند كتابة أي صيغة كمية تحديد: طول القطعة، نصف قطر الكرة، الوزن الناعم، والأس \(p\)، وعدم خلط الصيغ الموزونة بغير الموزونة.

\begin{remark}[استقلال المسارين]
\label{rem:ch25-independent-methods}
\resultid{ANT-PRIN-25-02}
\noindent\textenglish{METHODOLOGICAL-PRINCIPLE / INFERENCE-GUARDED}.
التوافق الفعّال يستخرج البنية من التطابقات ورفعها عبر قوى الأوليات، أما فك الاقتران فيفصل مساهمات قطع منحنى ذي هندسة محددة داخل فضاءات \(L^p\). اشتراكهما في النتيجة وفي وجود تكرار متعدد المقاييس لا يثبت تكافؤًا تقنيًا بينهما.
\end{remark}

\section{الجسر إلى وارينغ والمجاميع الأسية}
\symbolindex{الجسر إلى وارينغ والمجاميع الأسية}

\theoremindex{جسر منضبط إلى التطبيقات}
\begin{corollary}[جسر منضبط إلى التطبيقات]
\label{cor:ch25-waring-exponential-bridge}
\resultid{ANT-COR-25-01}
\noindent\textenglish{DERIVED-BRIDGE / INFERENCE-GUARDED}.
تمنح مبرهنة VMVT مدخلًا حاسمًا لتحسين تقديرات المجاميع الأسية متعددة الدرجات، كما تدخل في ضبط الأقواس الصغرى في مسائل وارينغ. لكن أي نتيجة كمية نهائية تحتاج، فوق VMVT، إلى اختيار متعدد الحدود، وتقريب ديوفانتي، وتجزئة الأقواس، وتقدير التكاملات الكبرى أو الصغرى بحسب المسألة.
\end{corollary}

وعليه لا نستنتج من المبرهنة وحدها صيغة تقاربية كاملة لعدد تمثيلات عدد صحيح كمجموع قوى. الفصل \ref{ch:circle-method-goldbach-waring} يوفر بنية الطريقة الدائرية، والفصل \ref{ch:exponential-sums-van-der-corput} يوفر أدوات تقدير المجاميع الأسية؛ أما الفصل الحالي فيضيف متوسطًا عاليًا أمثل يرفد تلك الأدوات.

\section{ما الذي لا تثبته المبرهنة؟}

\begin{itemize}
  \item لا تعطي VMVT تقديرًا نقطيًا أمثل لكل \(\boldsymbol\alpha\)؛ فهي مبرهنة متوسطية.
  \item لا تحول كل مسألة متعددة الحدود تلقائيًا إلى النظام القياسي \((n,n^2,\ldots,n^k)\).
  \item لا تستبدل التحليل المحلي في الأقواس الكبرى ولا فحص المتسلسلة المفردة في وارينغ.
  \item لا تجعل فك الاقتران والتوافق الفعّال قابلين للتبادل في جميع البيئات؛ لكل منهما مجال مرونة مختلف.
\end{itemize}

\section{الجبهة المفتوحة}
\symbolindex{الجبهة المفتوحة}

\begin{openproblem}[أنظمة متعددة الحدود وفك الاقتران الأوسع]
\label{open:ch25-polynomial-systems}
\resultid{ANT-OPEN-25-01}
\openresult
توسيع الصورة المثلى إلى أنظمة متعددة حدود عامة ذات ورونسكي غير منعدم، وإلى بيئات حقول الأعداد وحقول الدوال، وإلى مسائل فك اقتران لمتشعبات أوسع، يفتح عائلة كبيرة من الأسئلة. توجد نتائج قوية في هذه الاتجاهات، لكن هذا الفصل لا يدعي تغطيتها أو استخراج صيغة عامة موحدة منها.
\end{openproblem}

\section*{خلاصة الفصل}

حوّل التعامد قيمة متوسطة تحليلية إلى عد دقيق لحلول نظام فينوغرادوف. وكشف الحد القطري وصندوق الأطوار الصغيرة عن الحدين البنيويين \(X^s\) و\(X^{2s-K}\). أما مطابقة هذين الحدين من الأعلى فهي الإنجاز العميق الذي تحقق بمسارين مستقلين: التوافق الفعّال عند Wooley، وفك الاقتران للمنحنى اللحظي عند Bourgain--Demeter--Guth. قوة النتيجة لا تلغي الحاجة إلى الأدوات الخاصة بكل تطبيق، ولا تمنح إذنًا بخلط التاريخ أو التطبيعات أو أنواع البرهان.

\section*{تمارين}
\begin{enumerate}
  \item أثبت هوية التعامد في القضية \ref{prop:ch25-orthogonality-count} أولًا للحالة \(k=1\)، ثم عممها.
  \item احسب الوزن الحرج للحالات \(k=2,3,4\)، وحدد موضع الانتقال بين الحدين في كل حالة.
  \item أعط عددًا أدنى أدق للحلول القطرية يسمح بتبديل عناصر \((x_1,\ldots,x_s)\) عندما تكون متميزة.
  \item اكتب بالتفصيل اختيار \(c_k\) في برهان القضية \ref{prop:ch25-structural-lower-bounds} بحيث تحصل على ثابت صريح في \(|f_k|\gg_k X\).
  \item اشرح لماذا لا تكفي معرفة \(J_{s,k}(X)\) لاستخراج حد نقطي موحد لـ\(|f_k(\boldsymbol\alpha;X)|\).
  \item قارن بين مصدر تعدد المقاييس في التوافق الفعّال ومصدره في فك الاقتران، من دون ادعاء التكافؤ.
\end{enumerate}
